%% ============================================================
%%  SECCIÓN 3 — Principio Generalizado de Robertson-Schrödinger
%% ============================================================

\section{El Principio Generalizado de Incertidumbre de Robertson--\texorpdfstring{Schrödinger}{Schrodinger}}
\label{sec:incertidumbre}

\subsection{La Desigualdad de Cauchy--Schwarz en el Espacio de Hilbert}

El fundamento geométrico de todas las relaciones de indeterminación cuántica reside en la desigualdad de Cauchy--Schwarz, la cual constituye el análogo en espacios vectoriales complejos de la propiedad geométrica elemental $|\vec{a}\cdot\vec{b}| \le \norm{\vec{a}}\norm{\vec{b}}$ del espacio euclídeo \cite{zettili}.

\begin{teoremabox}[Teorema: Desigualdad de Cauchy--Schwarz en $\mathcal{H}$]
    Para cualesquiera dos vectores $\ket{u}, \ket{v} \in \mathcal{H}$, se cumple estrictamente:
    \begin{equation}
        |\braket{u}{v}|^2 \le \braket{u}{u}\braket{v}{v} = \norm{u}^2 \norm{v}^2,
        \label{eq:cauchy_schwarz}
    \end{equation}
    alcanzándose la igualdad si y solo si los vectores son linealmente dependientes (es decir, $\ket{u} = c\ket{v}$ para algún escalar $c \in \mathbb{C}$, o alguno de ellos es el vector nulo).
\end{teoremabox}

\textbf{Demostración Rigurosa Sin Atajos:}\\
Si $\ket{v} = 0$, la desigualdad se cumple trivialmente ($0 \le 0$). Supongamos entonces que $\ket{v} \neq 0$, lo que implica $\braket{v}{v} > 0$ por el axioma de positividad definida. Consideremos un vector auxiliar $\ket{\phi}$ construido como una combinación lineal con un parámetro complejo arbitrario $\lambda \in \mathbb{C}$:
\begin{equation}
    \ket{\phi} = \ket{u} + \lambda\ket{v}.
\end{equation}
Por el axioma de positividad definida del producto interno, la norma al cuadrado de $\ket{\phi}$ no puede ser negativa:
\begin{align}
    0 \le \braket{\phi}{\phi} &= \braket{u + \lambda v}{u + \lambda v} \nonumber\\
    &= \braket{u}{u} + \lambda\braket{u}{v} + \lambda^*\braket{v}{u} + \lambda^*\lambda\braket{v}{v} \nonumber\\
    &= \braket{u}{u} + \lambda\braket{u}{v} + \lambda^*\braket{u}{v}^* + |\lambda|^2\braket{v}{v}.
    \label{eq:cs_expansion}
\end{align}
Dado que esta desigualdad debe ser válida para \textit{cualquier} número complejo $\lambda$, elegimos el valor que minimiza la forma cuadrática, el cual corresponde a la proyección ortogonal de $\ket{u}$ sobre $\ket{v}$:
\begin{equation}
    \lambda = -\frac{\braket{v}{u}}{\braket{v}{v}}.
    \label{eq:optimal_lambda}
\end{equation}
Tomando el conjugado complejo de $\lambda$:
\begin{equation}
    \lambda^* = -\frac{\braket{v}{u}^*}{\braket{v}{v}^*} = -\frac{\braket{u}{v}}{\braket{v}{v}}.
\end{equation}
Sustituyendo cuidadosamente $\lambda$ y $\lambda^*$ en la Ec.~\eqref{eq:cs_expansion}:
\begin{align}
    0 &\le \braket{u}{u} + \left(-\frac{\braket{v}{u}}{\braket{v}{v}}\right)\braket{u}{v} + \left(-\frac{\braket{u}{v}}{\braket{v}{v}}\right)\braket{v}{u} + \left|\frac{\braket{v}{u}}{\braket{v}{v}}\right|^2 \braket{v}{v} \nonumber\\
    &= \braket{u}{u} - \frac{\braket{v}{u}\braket{u}{v}}{\braket{v}{v}} - \frac{\braket{u}{v}\braket{v}{u}}{\braket{v}{v}} + \frac{|\braket{v}{u}|^2}{\braket{v}{v}^2}\braket{v}{v}.
\end{align}
Recordando que $\braket{v}{u}\braket{u}{v} = \braket{u}{v}^*\braket{u}{v} = |\braket{u}{v}|^2$, los términos se simplifican algebraicamente:
\begin{equation}
    0 \le \braket{u}{u} - \frac{|\braket{u}{v}|^2}{\braket{v}{v}} - \frac{|\braket{u}{v}|^2}{\braket{v}{v}} + \frac{|\braket{u}{v}|^2}{\braket{v}{v}} = \braket{u}{u} - \frac{|\braket{u}{v}|^2}{\braket{v}{v}}.
\end{equation}
Multiplicando toda la desigualdad por el número real positivo $\braket{v}{v} > 0$:
\begin{equation}
    0 \le \braket{u}{u}\braket{v}{v} - |\braket{u}{v}|^2 \implies \boxed{|\braket{u}{v}|^2 \le \braket{u}{u}\braket{v}{v}}.
\end{equation}
Queda formalmente demostrada la desigualdad de Cauchy--Schwarz sin recurrir a derivadas de funciones reales ni asumir que el producto interno sea real \cite{sakurai}. Nótese que la igualdad en la Ec.~\eqref{eq:cauchy_schwarz} se alcanza si y solo si $\braket{\phi}{\phi} = 0 \iff \ket{\phi} = 0$, lo cual equivale a la dependencia lineal $\ket{u} = -\lambda\ket{v}$ entre ambos estados (con el escalar óptimo $\lambda = -\braket{v}{u}/\braket{v}{v}$).

\begin{figure}[H]
\centering
\begin{tikzpicture}[scale=1.15]
  \coordinate (O) at (0,0);
  \coordinate (U) at (3.5,2.5);
  \coordinate (V) at (5,1);
  \coordinate (P) at (3.846,0.769);
  \draw[-{Stealth},very thick,uniblue] (O) -- (V) node[right] {$\ket v$};
  \draw[-{Stealth},very thick,unired] (O) -- (U) node[above] {$\ket u$};
  \draw[-{Stealth},very thick,unigreen] (O) -- (P)
        node[midway,below,font=\small] {$\dfrac{\braket vu}{\braket vv}\ket v$};
  \draw[dashed,thick,black] (P) -- (U)
        node[midway,left=3pt,font=\small]
        {$\ket\phi$};
  \draw (4.140,0.828) -- (4.081,1.122) -- (3.787,1.063);
  \draw pic["$\vartheta$",draw=unigray,angle radius=1.1cm,angle eccentricity=1.35]
        {angle=V--O--U};
  \node[font=\small,align=left] at (7.2,3.0)
    {$\ket\phi=\ket u-\dfrac{\braket vu}{\braket vv}\ket v$\\[4pt]
     $\braket\phi v=0$ \\[2pt] $\Rightarrow\ \braket\phi\phi\ge0$ \\[2pt]
     $\Rightarrow\ |\braket uv|^2\le\braket uu\braket vv$};
\end{tikzpicture}
\caption{Cauchy--Schwarz: $\lambda=-\braket vu/\braket vv$ hace que $\ket\phi=\ket u+\lambda\ket v$ sea el residuo ortogonal a $\ket v$; su norma no puede ser negativa.}
\label{fig:cauchy_schwarz}
\end{figure}

\subsection{Construcción de los Estados de Desviación}

Sean $\hat{A}$ y $\hat{B}$ dos observables físicos representados por operadores lineales hermíticos ($\hat{A}\adj = \hat{A}$, $\hat{B}\adj = \hat{B}$). Supongamos que el sistema físico se encuentra en un estado normalizado arbitrario $\ket{\psi}$ ($\braket{\psi}{\psi}=1$).

Definimos los estados desviados (o fluctuaciones de estado) aplicando los operadores desviación definidos en la Ec.~\eqref{eq:deviation_operator}:
\begin{equation}
    \ket{u} \equiv \Delta\hat{A}\ket{\psi} = (\hat{A} - \ev{\hat{A}}\hat{I})\ket{\psi},
    \label{eq:state_u}
\end{equation}
\begin{equation}
    \ket{v} \equiv \Delta\hat{B}\ket{\psi} = (\hat{B} - \ev{\hat{B}}\hat{I})\ket{\psi}.
    \label{eq:state_v}
\end{equation}
Evaluemos ahora las normas al cuadrado de estos vectores:
\begin{equation}
    \braket{u}{u} = \braket{\Delta\hat{A}\psi}{\Delta\hat{A}\psi} = \matrixel{\psi}{(\Delta\hat{A})\adj\Delta\hat{A}}{\psi}.
\end{equation}
Debido a que $\Delta\hat{A}$ es hermítico ($(\Delta\hat{A})\adj = \Delta\hat{A}$):
\begin{equation}
    \braket{u}{u} = \matrixel{\psi}{(\Delta\hat{A})^2}{\psi} = (\Delta A)^2.
    \label{eq:norm_u_sq}
\end{equation}
Exactamente con el mismo procedimiento para el estado $\ket{v}$:
\begin{equation}
    \braket{v}{v} = \matrixel{\psi}{(\Delta\hat{B})^2}{\psi} = (\Delta B)^2.
    \label{eq:norm_v_sq}
\end{equation}
Por otro lado, el producto interno cruzado entre ambos estados desviados resulta ser:
\begin{equation}
    \braket{u}{v} = \matrixel{\psi}{(\Delta\hat{A})\adj\Delta\hat{B}}{\psi} = \matrixel{\psi}{\Delta\hat{A}\Delta\hat{B}}{\psi} = \ev{\Delta\hat{A}\Delta\hat{B}}.
    \label{eq:inner_uv}
\end{equation}
Sustituyendo directamente los resultados de las Ecs.~\eqref{eq:norm_u_sq}, \eqref{eq:norm_v_sq} y \eqref{eq:inner_uv} en la desigualdad de Cauchy--Schwarz \eqref{eq:cauchy_schwarz}:
\begin{equation}
    \boxed{|\ev{\Delta\hat{A}\Delta\hat{B}}|^2 \le (\Delta A)^2 (\Delta B)^2}.
    \label{eq:cs_applied}
\end{equation}

\subsection{Descomposición Algebraica: Conmutador y Anticonmutador}

El producto de dos operadores hermíticos arbitrarios $\Delta\hat{A}\Delta\hat{B}$ no es en general hermítico, ya que $(\Delta\hat{A}\Delta\hat{B})\adj = \Delta\hat{B}\adj\Delta\hat{A}\adj = \Delta\hat{B}\Delta\hat{A} \neq \Delta\hat{A}\Delta\hat{B}$ a menos que conmuten.

Sin embargo, cualquier producto de operadores puede descomponerse de forma exacta y única en una \textbf{parte simétrica} y una \textbf{parte antisimétrica} mediante la identidad:
\begin{equation}
    \Delta\hat{A}\Delta\hat{B} = \frac{1}{2}\{\Delta\hat{A}, \Delta\hat{B}\} + \frac{1}{2}[\Delta\hat{A}, \Delta\hat{B}],
    \label{eq:symmetric_antisymmetric_decomp}
\end{equation}
donde $\{\hat{X}, \hat{Y}\} \equiv \hat{X}\hat{Y} + \hat{Y}\hat{X}$ es el \textbf{anticonmutador} y $[\hat{X}, \hat{Y}] \equiv \hat{X}\hat{Y} - \hat{Y}\hat{X}$ es el \textbf{conmutador}.

\subsubsection{Cálculo Detallado del Conmutador de Desviaciones}
Desarrollemos algebraicamente el conmutador $[\Delta\hat{A}, \Delta\hat{B}]$:
\begin{align}
    [\Delta\hat{A}, \Delta\hat{B}] &= [\hat{A} - \ev{\hat{A}}\hat{I}, \, \hat{B} - \ev{\hat{B}}\hat{I}] \nonumber\\
    &= [\hat{A}, \hat{B}] - [\hat{A}, \ev{\hat{B}}\hat{I}] - [\ev{\hat{A}}\hat{I}, \hat{B}] + [\ev{\hat{A}}\hat{I}, \ev{\hat{B}}\hat{I}].
\end{align}
Dado que los valores esperados $\ev{\hat{A}}$ y $\ev{\hat{B}}$ son números escalares ordinarios (c-números) y la identidad $\hat{I}$ conmuta con absolutamente cualquier operador ($[\hat{X}, c\hat{I}] = c[\hat{X}, \hat{I}] = 0$):
\begin{equation}
    \boxed{[\Delta\hat{A}, \Delta\hat{B}] = [\hat{A}, \hat{B}]}.
    \label{eq:comm_dev_eq}
\end{equation}
¡El conmutador de las fluctuaciones es exactamente idéntico al conmutador de los operadores originales! Las medias estadísticas constantes se cancelan por completo.

\subsubsection{Propiedades de Hermiticidad de las Partes}
Analicemos la naturaleza de cada operador bajo la conjugación hermítica:
\begin{itemize}
    \item \textbf{El anticonmutador $\{\Delta\hat{A}, \Delta\hat{B}\}$ es HERMÍTICO:}
    \begin{align}
        \{\Delta\hat{A}, \Delta\hat{B}\}\adj &= (\Delta\hat{A}\Delta\hat{B} + \Delta\hat{B}\Delta\hat{A})\adj = \Delta\hat{B}\adj\Delta\hat{A}\adj + \Delta\hat{A}\adj\Delta\hat{B}\adj \nonumber\\
        &= \Delta\hat{B}\Delta\hat{A} + \Delta\hat{A}\Delta\hat{B} = \{\Delta\hat{A}, \Delta\hat{B}\}.
    \end{align}
    Por consiguiente, su valor esperado es obligatoriamente un \textbf{número real puro}:
    \begin{equation}
        \ev{\{\Delta\hat{A}, \Delta\hat{B}\}} \in \mathbb{R}.
    \end{equation}
    Esta cantidad real define la \textbf{covarianza cuántica}:
    \begin{equation}
        \mathrm{Cov}(\hat{A}, \hat{B}) \equiv \frac{1}{2}\ev{\{\Delta\hat{A}, \Delta\hat{B}\}} = \frac{1}{2}\ev{\hat{A}\hat{B} + \hat{B}\hat{A}} - \ev{\hat{A}}\ev{\hat{B}}.
    \end{equation}
    
    \item \textbf{El conmutador $[\hat{A}, \hat{B}]$ es ANTIHERMÍTICO:}
    \begin{equation}
        [\hat{A}, \hat{B}]\adj = (\hat{A}\hat{B} - \hat{B}\hat{A})\adj = \hat{B}\adj\hat{A}\adj - \hat{A}\adj\hat{B}\adj = \hat{B}\hat{A} - \hat{A}\hat{B} = -[\hat{A}, \hat{B}].
    \end{equation}
    Por consiguiente, su valor esperado es estrictamente un \textbf{número imaginario puro}:
    \begin{equation}
        \ev{[\hat{A}, \hat{B}]} \in \iu\mathbb{R} \implies \ev{[\hat{A}, \hat{B}]}^* = -\ev{[\hat{A}, \hat{B}]}.
    \end{equation}
\end{itemize}

\subsection{Deducción de la Desigualdad de Robertson--\texorpdfstring{Schrödinger}{Schrodinger}}

Tomando el valor esperado de la descomposición \eqref{eq:symmetric_antisymmetric_decomp}:
\begin{equation}
    \ev{\Delta\hat{A}\Delta\hat{B}} = \underbrace{\frac{1}{2}\ev{\{\Delta\hat{A}, \Delta\hat{B}\}}}_{\text{Parte Real } X} + \underbrace{\frac{1}{2}\ev{[\hat{A}, \hat{B}]}}_{\text{Parte Imaginaria } \iu Y}.
\end{equation}
Para cualquier número complejo $z = X + \iu Y$ con $X, Y \in \mathbb{R}$, el cuadrado de su módulo se calcula de forma exacta como la suma de los cuadrados de sus partes real e imaginaria:
\begin{equation}
    |z|^2 = (\mathrm{Re}\,z)^2 + (\mathrm{Im}\,z)^2 = X^2 + Y^2.
\end{equation}
Identificando los términos:
\begin{align}
    \mathrm{Re}\left(\ev{\Delta\hat{A}\Delta\hat{B}}\right) &= \frac{1}{2}\ev{\{\Delta\hat{A}, \Delta\hat{B}\}}, \\
    \mathrm{Im}\left(\ev{\Delta\hat{A}\Delta\hat{B}}\right) &= \frac{1}{2\iu}\ev{[\hat{A}, \hat{B}]} = -\frac{\iu}{2}\ev{[\hat{A}, \hat{B}]}.
\end{align}
Por lo tanto:
\begin{equation}
    |\ev{\Delta\hat{A}\Delta\hat{B}}|^2 = \frac{1}{4}\left|\ev{\{\Delta\hat{A}, \Delta\hat{B}\}}\right|^2 + \frac{1}{4}\left|\ev{[\hat{A}, \hat{B}]}\right|^2.
    \label{eq:modulus_expanded}
\end{equation}
Sustituyendo la Ec.~\eqref{eq:modulus_expanded} en la desigualdad de Cauchy--Schwarz \eqref{eq:cs_applied}, obtenemos el teorema fundamental:

\begin{teoremabox}[Principio Generalizado de Incertidumbre de Robertson--Schrödinger (1930)]
    Para dos observables hermíticos arbitrarios $\hat{A}$ y $\hat{B}$ medidos sobre un estado cuántico $\ket{\psi}$:
    \begin{equation}
        \boxed{(\Delta A)^2 (\Delta B)^2 \ge \frac{1}{4}\left|\ev{\{\Delta\hat{A}, \Delta\hat{B}\}}\right|^2 + \frac{1}{4}\left|\ev{[\hat{A}, \hat{B}]}\right|^2}.
        \label{eq:robertson_schrodinger}
    \end{equation}
\end{teoremabox}

\subsection{Reducción a la Relación de Incertidumbre de Heisenberg}

\begin{alertabox}[¿Por qué puede descartarse el término del anticonmutador?]
    El término que contiene el anticonmutador es el cuadrado de un número real:
    \begin{equation}
        \frac{1}{4}\left|\ev{\{\Delta\hat{A}, \Delta\hat{B}\}}\right|^2 \ge 0.
    \end{equation}
    Al tratarse de una cantidad estrictamente no negativa, eliminarla del miembro derecho de la desigualdad produce una cota inferior más débil, pero igualmente válida:
    \begin{equation}
        (\Delta A)^2 (\Delta B)^2 \ge \frac{1}{4}\left|\ev{[\hat{A}, \hat{B}]}\right|^2.
    \end{equation}
    Extrayendo la raíz cuadrada a ambos lados, obtenemos la famosísima \textbf{Relación de Incertidumbre de Heisenberg--Robertson (1929)}:
    \begin{equation}
        \boxed{\Delta A \, \Delta B \ge \frac{1}{2}\left|\ev{[\hat{A}, \hat{B}]}\right|}.
        \label{eq:heisenberg_robertson}
    \end{equation}
    \textbf{Advertencia Física:} Si bien la relación de Heisenberg \eqref{eq:heisenberg_robertson} es la más conocida en los textos, la desigualdad completa de Robertson--Schrödinger \eqref{eq:robertson_schrodinger} es más general y siempre al menos tan restrictiva como la de Heisenberg (siendo estrictamente más restrictiva siempre que la covarianza cuántica $\mathrm{Cov}(\hat{A}, \hat{B}) \neq 0$). Existen estados cuánticos correlacionados (como los \textit{estados comprimidos} o \textit{squeezed states} en óptica cuántica) donde la covarianza cuántica $\mathrm{Cov}(\hat{A}, \hat{B}) \neq 0$ juega un papel crucial, haciendo que el producto $\Delta A \Delta B$ sea estrictamente mayor que $\frac{1}{2}|\ev{[\hat{A}, \hat{B}]}|$ \cite{ballentine}.
\end{alertabox}

\subsection{Ejemplo Resuelto 2: Relación de Incertidumbre Posición--Momento y Paquete Gaussiano}

\begin{ejemplobox}[Ejemplo Resuelto 2: Incertidumbre Mínima y Paquete de Ondas Coherente]
    \textbf{Enunciado:}
    \begin{enumerate}
        \item Utilice la relación de Heisenberg--Robertson \eqref{eq:heisenberg_robertson} para hallar el límite inferior del producto de incertidumbres de la posición $\hat{x}$ y el momento lineal $\hat{p}_x$, sabiendo que el conmutador canónico es $[\hat{x}, \hat{p}_x] = \iu\hbar\hat{I}$.
        \item Considere una partícula libre descrita por una función de onda gaussiana real normalizada:
        \begin{equation}
            \psi(x) = \left(\frac{1}{\pi \sigma^2}\right)^{1/4} \exp\left(-\frac{x^2}{2\sigma^2}\right).
        \end{equation}
        Calcule $\ev{\hat{x}}$, $\ev{\hat{p}_x}$, $\Delta x$, $\Delta p_x$ y evalúe la covarianza cuántica $\mathrm{Cov}(\hat{x}, \hat{p}_x) = \frac{1}{2}\ev{\{\Delta\hat{x}, \Delta\hat{p}_x\}}$.
    \end{enumerate}

    \textbf{Solución Paso a Paso:}

    \textit{Parte 1: Cota Teórica:}
    Evaluamos el valor esperado del conmutador canónico:
    \begin{equation}
        \ev{[\hat{x}, \hat{p}_x]} = \ev{\iu\hbar\hat{I}} = \iu\hbar\braket{\psi}{\psi} = \iu\hbar.
    \end{equation}
    El módulo es $|\iu\hbar| = \hbar$. Aplicando directamente la Ec.~\eqref{eq:heisenberg_robertson}:
    \begin{equation}
        \Delta x \, \Delta p_x \ge \frac{1}{2}|\iu\hbar| = \frac{\hbar}{2}.
    \end{equation}
    Esta es la cota universal de Heisenberg para coordenadas cartesianas conjugadas.

    \textit{Parte 2: Cálculo Explícito en el Paquete Gaussiano:}
    Dado que $|\psi(x)|^2$ es una función par en $x \in (-\infty, \infty)$:
    \begin{equation}
        \ev{\hat{x}} = \int_{-\infty}^\infty x |\psi(x)|^2\,\dif x = 0.
    \end{equation}
    Calculamos $\ev{\hat{x}^2}$ mediante la integral gaussiana $\int_{-\infty}^\infty x^2 \eu^{-\alpha x^2}\,\dif x = \frac{1}{2}\sqrt{\frac{\pi}{\alpha^3}}$ con $\alpha = 1/\sigma^2$:
    \begin{equation}
        \ev{\hat{x}^2} = \left(\frac{1}{\pi\sigma^2}\right)^{1/2} \int_{-\infty}^\infty x^2 \eu^{-x^2/\sigma^2}\,\dif x = \frac{1}{\sqrt{\pi}\sigma} \cdot \frac{\sqrt{\pi}}{2}\sigma^3 = \frac{\sigma^2}{2}.
    \end{equation}
    Por ende:
    \begin{equation}
        \Delta x = \sqrt{\ev{\hat{x}^2} - \ev{\hat{x}}^2} = \sqrt{\frac{\sigma^2}{2}} = \frac{\sigma}{\sqrt{2}}.
    \end{equation}
    Para el momento lineal, actuamos con el operador diferencial $\hat{p}_x = -\iu\hbar\,\frac{\dif}{\dif x}$:
    \begin{equation}
        \hat{p}_x \psi(x) = -\iu\hbar \frac{\dif}{\dif x}\left[ C \eu^{-x^2/(2\sigma^2)} \right] = -\iu\hbar \left(-\frac{x}{\sigma^2}\right)\psi(x) = \iu\hbar\frac{x}{\sigma^2}\psi(x).
    \end{equation}
    El valor esperado del momento es nulo por antisimetría:
    \begin{equation}
        \ev{\hat{p}_x} = \int_{-\infty}^\infty \psi^*(x)\hat{p}_x\psi(x)\,\dif x = \frac{\iu\hbar}{\sigma^2}\int_{-\infty}^\infty x |\psi(x)|^2\,\dif x = 0.
    \end{equation}
    Para $\ev{\hat{p}_x^2}$, usamos la propiedad de que $\hat{p}_x$ es hermítico ($\braket{\psi}{\hat{p}_x^2\psi} = \braket{\hat{p}_x\psi}{\hat{p}_x\psi}$):
    \begin{align}
        \ev{\hat{p}_x^2} &= \int_{-\infty}^\infty |\hat{p}_x\psi(x)|^2\,\dif x = \int_{-\infty}^\infty \left|\iu\hbar\frac{x}{\sigma^2}\psi(x)\right|^2\,\dif x \nonumber\\
        &= \frac{\hbar^2}{\sigma^4}\int_{-\infty}^\infty x^2 |\psi(x)|^2\,\dif x = \frac{\hbar^2}{\sigma^4}\ev{\hat{x}^2} = \frac{\hbar^2}{\sigma^4}\left(\frac{\sigma^2}{2}\right) = \frac{\hbar^2}{2\sigma^2}.
    \end{align}
    Por ende:
    \begin{equation}
        \Delta p_x = \sqrt{\ev{\hat{p}_x^2} - \ev{\hat{p}_x}^2} = \frac{\hbar}{\sqrt{2}\sigma}.
    \end{equation}
    Multiplicando las dos incertidumbres:
    \begin{equation}
        \Delta x \, \Delta p_x = \left(\frac{\sigma}{\sqrt{2}}\right)\left(\frac{\hbar}{\sqrt{2}\sigma}\right) = \frac{\hbar}{2}.
    \end{equation}
    ¡El paquete gaussiano satura con precisión absoluta la igualdad de Heisenberg!

    Evaluemos finalmente la covarianza: como $\ev{\hat{x}} = \ev{\hat{p}_x} = 0$, se tiene $\Delta\hat{x} = \hat{x}$ y $\Delta\hat{p}_x = \hat{p}_x$. Calculando el valor esperado del producto ordenado:
    \begin{align}
        \ev{\hat{x}\hat{p}_x} &= \int_{-\infty}^\infty \psi^*(x)\hat{x}\left(\iu\hbar\frac{x}{\sigma^2}\psi(x)\right)\dif x = \frac{\iu\hbar}{\sigma^2}\int_{-\infty}^\infty x^2 |\psi(x)|^2\,\dif x \nonumber\\
        &= \frac{\iu\hbar}{\sigma^2}\ev{\hat{x}^2} = \frac{\iu\hbar}{\sigma^2}\left(\frac{\sigma^2}{2}\right) = \frac{\iu\hbar}{2}.
        \label{a2-eq:xp_exp}
    \end{align}
    Puesto que $\hat{x}$ y $\hat{p}_x$ son operadores hermíticos, $(\hat{x}\hat{p}_x)\adj = \hat{p}_x\hat{x}$, de modo que el valor esperado del producto inverso es su complejo conjugado:
    \begin{equation}
        \ev{\hat{p}_x\hat{x}} = \ev{\hat{x}\hat{p}_x}^* = \left(\frac{\iu\hbar}{2}\right)^* = -\frac{\iu\hbar}{2}.
        \label{a2-eq:px_exp}
    \end{equation}
    Por consiguiente, el valor esperado del anticonmutador se anula de forma exacta:
    \begin{equation}
        \ev{\{\hat{x}, \hat{p}_x\}} = \ev{\hat{x}\hat{p}_x + \hat{p}_x\hat{x}} = \ev{\hat{x}\hat{p}_x} + \ev{\hat{p}_x\hat{x}} = \frac{\iu\hbar}{2} - \frac{\iu\hbar}{2} = 0 \implies \mathrm{Cov}(\hat{x}, \hat{p}_x) = 0.
    \end{equation}
    Esto demuestra por cálculo explícito que para el paquete gaussiano mínimo la covarianza cuántica es estrictamente nula, de modo que las cotas de Heisenberg y Robertson--Schrödinger coinciden en $\hbar/2$.
\end{ejemplobox}

\begin{figure}[H]
    \centering
    % Fig. 1: paquete gaussiano en x y en p (hbar = 1)
\begin{tikzpicture}
\begin{groupplot}[
  group style={group size=2 by 1, horizontal sep=1.8cm},
  width=0.48\linewidth, height=5.2cm, domain=-4:4, samples=200,
  axis lines=left, xlabel near ticks, ylabel near ticks,
  legend style={font=\scriptsize, draw=none, fill=none},
  every axis plot/.append style={thick}]
\nextgroupplot[xlabel={$x$ (unidades arbitrarias de longitud, $\hbar=1$)}, ylabel={$|\psi(x)|^2$},
  xlabel style={font=\scriptsize}, title={\small Espacio de posiciones}, ymin=0, ymax=1.35,
  legend pos=north east]
  \addplot[uni-blue]{exp(-x^2/0.5^2)/(sqrt(pi)*0.5)};
  \addplot[alert-border]{exp(-x^2/1.5^2)/(sqrt(pi)*1.5)};
  \legend{$\sigma=0.5$, $\sigma=1.5$}
\nextgroupplot[xlabel={$p$ (unidades arbitrarias de momento, $\hbar=1$)}, ylabel={$|\tilde\psi(p)|^2$},
  xlabel style={font=\scriptsize}, title={\small Espacio de momentos}, ymin=0, ymax=1.35, legend pos=north east]
  \addplot[uni-blue]{0.5*exp(-0.5^2*x^2)/sqrt(pi)};
  \addplot[alert-border]{1.5*exp(-1.5^2*x^2)/sqrt(pi)};
  \legend{$\sigma=0.5$, $\sigma=1.5$}
\end{groupplot}
\end{tikzpicture}

    \caption{Paquete gaussiano del Ejemplo Resuelto 2 en posici\'on y en momento ($\hbar=1$). Un $\sigma$ peque\~no localiza la part\'icula en $x$ pero ensancha $|\tilde\psi(p)|^2$, y viceversa; en ambos casos $\Delta x\,\Delta p_x=\hbar/2$.}
    \label{fig:gauss_xp}
\end{figure}
